% !PRINTED 414
\weberpage{414}
If, on the other hand, equation (1) is not an identity, then $\omega$ is
the root of a quadratic equation with integral coefficients:
\weq{2}{b\,\omega^2 + (a-\partial)\,\omega - c = 0,}
hence, if we put
\[
a\,\partial - b\,c = n
\]
\[
m = -4\,b\,c - (a-\partial)^2 = 4\,n - (a+\partial)^2
\]
then we have
\[
m \equiv 0 \ \text{ or } \ \equiv -1 \ (\mathrm{mod}\ 4),
\]
and we obtain:
\[
\omega = \frac{-a+\partial+\sqrt{-m}}{2\,b},
\]
\[
\mu = \frac{a+\partial+\sqrt{-m}}{2},
\]
\[
n = \frac{a+\partial+\sqrt{-m}}{2}\cdot\frac{a+\partial-\sqrt{-m}}{2}, % Source scan shows a minus; the mark above it is a scan/print speck.
\]
\[
\mu^2 - (a+\partial)\,\mu + n = 0.
\]

In order that $\omega$ may have a positive imaginary part, $m$, and all
the more therefore $n$, must be \gesperrt{positive}, and $\mu$ is a complex,
integral, algebraic number. This explains the designation
\gesperrt{complex multiplication}.

If now, conversely, $\omega$ satisfies a quadratic equation
\weq{3}{A\,\omega^2 + B\,\omega + C = 0}
in which $A$, $B$, $C$ are integers, and the discriminant
\weq{4}{D = B^2 - 4\,A\,C}
is negative, then $\omega$ is not real, and in one of the two roots
of (3) the imaginary part is positive. This one shall be taken for
$\omega$. Then $\omega$ is called a root of the quadratic form
$(A, B, C)$.

Without prejudice to generality, $A$, $B$, $C$ may be assumed to be
without a common divisor and $A$, $C$ positive. Infinitely many
relations of the form (1) [or (2)] can then be set up for $\omega$. One
has only to put, when $x$ is an integer different from zero:
\weq{5}{\begin{aligned}
b &= A\,x,\\
c &= -\,C\,x,\\
a-\partial &= B\,x.
\end{aligned}}
